% ==================================================================================================================================
% Introduction

\minitoc  % Affiche la table des matières pour ce chapitre

Dans ce chapitre, nous allons définir formellement les matrices ainsi que les espaces vectoriels de matrices. 
Ces objets vont nous permettre de représenter très facilement les applications linéaires introduites dans le 
chapitre 1. Les propriétés que nous allons voir sur les matrices pourront, pour la plupart se transposer aux applications 
représentées. 

\vspace{0.3cm}

Dans tout ce chapitre, on nommera $E$ un $\K$-espace vectoriel où $\K$ est un corps commutatif. 

% ==================================================================================================================================
% Généralités 

\section{Généralités}

\begin{definition}[Matrice]
    On appelle matrice $n \in \N$ lignes et $p \in \N$ colonnes toute application 
        \[ M : \{1, \dots , n\} \times \{1, \dots , p\} \longrightarrow \K \] 
    On dit alors que $M$ est de taille $n \times p$.
    On note $ \mathcal{M}_{n,p}(\K)$ l'ensemble des matrices de taille $n \times p $ à coefficients dans un corps $\K$.  
\end{definition}

\begin{example}
    Les matrices sont représentées de façon tabulaire. 
    Par exemple une matrice $M$ à 2 lignes et 3 colonnes à coefficients dans $\R$ peut être : 
        \[ M = 
            \begin{pmatrix}
                1 & 26 & \frac{3}{2} \\ 
                \pi & 6 & \sqrt{2} 
            \end{pmatrix} \] 
    On notera $M(i,j)$ ou $M_{i,j}$ le coefficient de $M$ de la i-ème ligne et j-ème colonne. 
    Ici, on a donc $M(1,1) = 1$ et $M(2, 3) = \sqrt{2}$. 
\end{example}

\begin{proposition}[Matrices Remarquables]
    Le nombre de lignes et de colonnes d'une matrice pouvant varier dans $\N \times \N$, il existe des 
    matrices dites remarquables : 
    \begin{itemize}
        \item \textbf{Matrice ligne} : une matrice de taille $ 1 \times p$ de la forme : $(x_1, \dots, x_p)$ 
        \item \textbf{Matrice colonne} : de la forme 
        $ \begin{pmatrix}
            x_1 \\ 
            \vdots \\ 
            x_p 
        \end{pmatrix}$ 
    \end{itemize} 
\end{proposition}

\begin{proposition}[Matrice Égales]
    On dit que deux matrices M et N $ \in \mathcal{M}_{n,m}(\K)$ sont égales ssi elles sont de même taille et 
    tout leurs coefficients sont égaux. 
        \[ \text{i.e} \quad \forall i,j \in \llbracket 1, n \rrbracket \times \llbracket 1, m \rrbracket, M(i,j) = N(i,j) \]  
\end{proposition}

\begin{definition}[Transposition]
    Soit $ \mathcal{M}_{n,m}(\K)$ l'ensemble des matrices à coefficients dans un corps $\K$. 
    On défini l'application transposition comme : 
        \[ T : 
            \begin{cases}
                \mathcal{M}_{n,m}(\K) \longrightarrow \mathcal{M}_{m,n}(\K) \\ 
                M = (m_{i,j}) \longrightarrow ^t M = (m_{j,i})
            \end{cases} \]
    Cette application prend une matrice $ n \times m$ et la retourne en une matrice $ m \times n $ en "transformant"
    ses lignes en colonnes. 
    Elle est \textbf{involutive}, autrement dit $ \forall M \in \mathcal{M}_{n,m}(\K) \; ^t (^t M) = M$ 
    et \textbf{linéaire}. 
\end{definition}

\begin{definition}[Matrice symétrique]
    On dit qu'une matrice $M \in \mathcal{M}_{n}(\K)$ est symétrique ssi $^tM = M$. 
\end{definition}


% ==================================================================================================================================
% Opérations et Structure 

\section{Opérations et Structure}

Dans cette section, nous considérons l'ensemble des matrices $ \mathcal{M}_{n,m}(\K)$ de taille $ n \times m $
à coefficients dans un corps $\K$. 

\subsection{Addition et Produit Externe}

\begin{definition}[Somme Matricielle]
    Soient $A,B \in \mathcal{M}_{n,p}(\K)$ deux matrices de même taille. 
    On définit la matrice $A + B \in \mathcal{M}_{n,p}(\K)$ comme la matrice dont chaque élément est 
    l'exacte somme des éléments de $A$ et $B$ situés à la même place. 
    Plus formellement, si $ \forall (i,j) \in \llbracket 1, n \rrbracket \times  \llbracket 1, p \rrbracket$, 
    $ A = (a_{i,j})$ et $ B = (b_{i,j})$ alors : 
        \[ A + B = (a_{i,j} + b_{i,j}) \]
\end{definition}

\begin{example}[Somme Matricielle]
    Soit $A,B \in \mathcal{M}_{3,2}(\R)$ de la forme : 
        \[ A = 
            \begin{pmatrix}
                1 & 3 \\ 
                6 & 9 \\ 
                2 & 0 \\ 
            \end{pmatrix}
            \quad B = 
            \begin{pmatrix}
                1/2 & 0 \\ 
                6 & 3 \\ 
                11 & 2
            \end{pmatrix}
            \quad \text{alors} \quad A + B = 
            \begin{pmatrix}
                3/2 & 3 \\ 
                12 & 12 \\ 
                13 & 2 
            \end{pmatrix} \] 
\end{example}

L'ensemble $ \mathcal{M}_{n,p}(\K)$ muni de l'addition comme nous venons de la définir forme un groupe abélien 
$(\mathcal{M}_{n,p}(\K), +)$. 

\begin{definition}[Produit externe]
    Soient $A \in \mathcal{M}_{n,p}(\K)$ et $ \lambda \in \K$ un scalaire. 
    On définit la matrice $\lambda . M$ comme le \emph{produit externe} de $M$ par le scalaire $\lambda$. 
    C'est une matrice de $\mathcal{M}_{n,p}(\K)$. 
    Elle est obtenue en multipliant chaque valeur de $M$ par le scalaire $\lambda$. 
    On a donc : 
        \[ \forall (i,j) \in \llbracket 1, n \rrbracket \times  \llbracket 1, p \rrbracket, \quad \text{tq} \quad M = (m_{i,j}) \quad \text{et} \quad \lambda.M = (\lambda \times m_{i,j}) \] 
\end{definition}

\begin{prop}[Opérations et transposition]
    Soient deux matrices $A,B \in \mathcal{M}_{n,p}(\K)$ et $ \lambda \in \K$. On a les propriétés suivantes : 
    \begin{itemize}
        \item $^t(A + B) = \; ^tA + \; ^tB $
        \item $^t(\lambda A) = \lambda \; ^t A $ 
    \end{itemize}
\end{prop}

\begin{definition}[Matrice Antisyméytrique]
    Nous avons définit les matrices symétriques comme les matrices involutives par transposition. 
    Les matrices \emph{antisymétriques}, sont de la forme $M \in \mathcal{M}_{n,p}(\K)$ telle que $^tM = - M$. 
\end{definition}


\subsection{Structure}

Ces opérations nous permettent de donner une structure à l'espace $ \mathcal{M}_{n,p}(\K)$. 

\begin{prop}[Structure]
    L'ensemble $\mathcal{M}_{n,p}(\K)$ muni des opérations usuelles (addition et multiplication par un scalaire) 
    possède une structure de $\K$-espace vectoriel. 
    
    Le vecteur nul $0_{\mathcal{M}_{n,p}(\K)}$ correspond à la matrice nulle $0_{n,p}$. 
    De même, l'opposé de la matrice $A = (a_{i,j})$ est la matrice $-A = (-a_{i,j})$. 
\end{prop}

